\section{Implications, Reproducibility, and Validity}
\label{sec:implications-validity}

\subsection{Implications for Fraud Benchmarking}

The results favor contract-indexed reporting over a pooled leaderboard. A model name and dataset name do not identify a result unless temporal masks, graph visibility, construction, selection rule, decision budget, and resource envelope are also fixed. This matters even in a small model grid: Elliptic changes its AUPRC leader under a graph-only intervention, while DGraphFin keeps the leader but changes its margin. On IBM AML-Data, AUPRC and F1 choose different configurations in 11 of 16 feasible contexts.

Three reporting rules follow. First, graph visibility should be treated as an experimental variable, with a feature-only control when possible. Second, ranking, thresholded decisions, review capacity, and calibration should be reported as distinct claim types. Third, an out-of-memory cell belongs in a feasibility table, not at the bottom of a performance table. These rules do not prescribe one deployment contract; they make the chosen contract inspectable.

Cross-dataset aggregation would be especially misleading here. Elliptic predicts illicit transaction nodes, DGraphFin predicts fraudulent account nodes, and IBM AML-Data predicts laundering transaction edges. Their positive rates and temporal units differ, and IBM is synthetic. We therefore use each dataset to test a scoped proposition rather than average them into a single benchmark score.

\subsection{Artifact and Reproduction Paths}

The artifact connects each paper number to a source file, filter, seed set, aggregation rule, evidence lock, and output location. It includes dataset/protocol cards, a typed claim ledger, prediction-manifest references, resource and exclusion records, deterministic analysis code, figure/table builders, and the support validator. This extends documentation practices such as datasheets, model cards, and benchmark cards by making the claim-to-evidence relation executable \cite{gebru2021datasheets,mitchell2019modelcards,sokol2025benchmarkcards}.

The primary IBM package contains 840 validated result files and 840 corresponding prediction exports across the measured Small/Medium cells. That count describes provenance coverage, not 840 independent configurations. The real-data visibility family contains 180 result rows and 180 prediction files. Canonical locks distinguish these from partial imports, compatibility aliases, diagnostics, and resource-blocked plans.

Two reproduction paths are provided. The no-training path validates locks and manifests, recomputes all statistics from saved results and predictions, rebuilds tables and figures, executes the support mutations, and compiles both documents. The full path reacquires each dataset under its own access terms and reruns the training families on suitable compute. Raw restricted datasets, credentials, local absolute paths, failed temporary archives, and platform metadata are excluded from the release. The package claims no external artifact badge.

\subsection{Threats to Validity}

\textbf{Internal validity.} The strict and isolated real-data protocols share training code and masks, but a software error could still alter visibility in an unintended way. The MLP invariance check, prediction manifests, paired seeds, and graph-mode records reduce this risk. Selection differs across experiment families: the real-data grid saves the best validation-F1 checkpoint, IBM uses fixed training schedules and a 0.5 threshold, and GraphSafe fits its gate on validation reliability. Comparisons are confined to a family for this reason.

The IBM early-to-late classifier uses labels from the first 50\% of transactions, but its shared node-history matrix is built from the first 60\%. Thus node features include label-free covariates from the 50--60\% interval, which is the first half of that protocol's validation window. This is transductive covariate access, not test-label leakage. We state it as part of the contract and do not call the early-to-late features ``first-50\%-only.'' The temporal-stress duplicate-label defect demonstrates why manifest completeness alone is insufficient; construct arguments must reach the executed harness.

\textbf{Construct validity.} AUPRC measures score ranking under class imbalance, F1 depends on a threshold, and AUROC can look strong when practical precision remains low. Review-budget recall is closer to an investigation queue but does not measure case resolution, deterrence, or financial loss. Equation~\eqref{eq:cost-risk} uses a hypothetical 5:1 cost ratio. The IBM labels describe synthetic laundering scenarios, whereas Elliptic and DGraphFin use different real-data entities and label processes. None is a complete proxy for a private bank's AML workflow.

\textbf{External validity.} Validated predictive evidence covers two real public graphs and four Small/Medium variants of a synthetic generator. It does not cover a private-bank deployment, IBM Large, or Medium GINE. The T4 memory boundary may move with sampling, implementation, or hardware. The tested architectures and constructions are informative controls, not an exhaustive survey of temporal or heterogeneous GNNs. Dataset access and licensing also limit exact independent reruns.

\textbf{Statistical conclusion validity.} Ten seeds permit paired analysis but still leave uncertainty, particularly on DGraphFin and the low-intensity IBM variants. Very small absolute AUPRC changes can have large relative values at prevalence below 0.2\%. Bootstrap intervals describe seed/cell variability in the observed grid, not sampling from all institutions or time periods. Deterministic library baselines can repeat exactly across nominal seeds; their ten rows document repeated pipeline evaluation, not ten stochastic fits. We use Holm correction for declared analysis families, keep descriptive rank correlations separate, and never pool unequal feasible sets. The four fixed IBM contexts and the six GraphSafe contexts are averaged within seed before inference; Small-only GINE means are not compared as though Medium were observed. Context-specific IBM rows remain visible because blocking can hide heterogeneity.

\subsection{Ethics and Intended Use}

Fraud flags can impose investigation costs, delay legitimate transactions, and expose people or organizations to adverse action. False negatives also carry social and financial costs. The reported review budgets are intended for human-in-the-loop triage, not automatic punitive decisions. The study does not evaluate demographic fairness because the validated artifacts do not contain an appropriate fairness surface; absence of such an analysis is not evidence of fairness. Synthetic IBM data reduces direct privacy exposure but cannot establish realism. Any operational use requires institution-specific validation, governance, appeal procedures, security review, and lawful data handling.
